\textbf{Phần 1. Trắc nghiệm lựa chọn}
\setcounterpageref{ex}{0}

\begin{ex} %[5D4V22]
Tìm $x$, biết: $x \times 4{,}5 = 31{,}5 - 6{,}75$.
\choice
{$4{,}5$}
{\True $5{,}5$}
{$6{,}5$}
{$5{,}25$}
\loigiai{
Ta có:
\begin{align*}
x \times 4{,}5 &= 24{,}75\\
x &= 24{,}75 : 4{,}5\\
x &= 5{,}5
\end{align*}
\textbf{Đáp án của câu hỏi là B}
}
\end{ex}

\begin{ex} %[5D4V22]
May mỗi bộ quần áo cần $2{,}8\text{ m}$ vải. Có một cuộn vải dài $65\text{ m}$ thì may được nhiều nhất bao nhiêu bộ quần áo như thế và còn thừa bao nhiêu mét vải?
\choice
{$22\text{ bộ và thừa }3{,}4\text{ m}$}
{\True $23\text{ bộ và thừa }0{,}6\text{ m}$}
{$23\text{ bộ và thừa }6\text{ m}$}
{$24\text{ bộ và thừa }0{,}2\text{ m}$}
\loigiai{
Thực hiện phép chia: $65 : 2{,}8 = 23$ (dư $0{,}6$).\\
Vậy may được nhiều nhất $23$ bộ quần áo và còn thừa $0{,}6\text{ m}$ vải.\\
\textbf{Đáp án của câu hỏi là B}
}
\end{ex}

\begin{ex} %[5D4V22]
Tính giá trị của biểu thức: $48 : 0{,}25 + 16 : 0{,}125$.
\choice
{$160$}
{$240$}
{\True $320$}
{$400$}
\loigiai{
Chia cho $0{,}25$ là nhân với $4$; chia cho $0{,}125$ là nhân với $8$:\\
$$48 \times 4 + 16 \times 8 = 192 + 128 = 320.$$
\textbf{Đáp án của câu hỏi là C}
}
\end{ex}

\begin{ex} %[5D4V22]
Một ô tô đi quãng đường $144\text{ km}$ hết $3{,}2$ giờ. Một xe máy cũng đi quãng đường đó hết $4{,}5$ giờ. Hỏi mỗi giờ ô tô đi nhiều hơn xe máy bao nhiêu ki-lô-mét?
\choice
{$11\text{ km}$}
{$12\text{ km}$}
{\True $13\text{ km}$}
{$15\text{ km}$}
\loigiai{
Mỗi giờ ô tô đi được số ki-lô-mét là:
$144 : 3{,}2 = 45\text{ (km)}$.\\
Mỗi giờ xe máy đi được số ki-lô-mét là:
$144 : 4{,}5 = 32\text{ (km)}$.\\
Mỗi giờ ô tô đi nhiều hơn xe máy là:
$45 - 32 = 13\text{ (km)}$.\\
\textbf{Đáp án của câu hỏi là C}
}
\end{ex}

\begin{ex} %[5D4V22]
Biết $4{,}5\text{ l}$ dầu hoả cân nặng $3{,}6\text{ kg}$. Hỏi $7\text{ l}$ dầu hoả đó cân nặng bao nhiêu ki-lô-gam?
\choice
{$5{,}2\text{ kg}$}
{\True $5{,}6\text{ kg}$}
{$6\text{ kg}$}
{$6{,}4\text{ kg}$}
\loigiai{
Mỗi lít dầu hoả cân nặng số ki-lô-gam là:
$3{,}6 : 4{,}5 = 0{,}8\text{ (kg)}$.\\
$7\text{ l}$ dầu hoả cân nặng số ki-lô-gam là:
$0{,}8 \times 7 = 5{,}6\text{ (kg)}$.\\
\textbf{Đáp án của câu hỏi là B}
}
\end{ex}

\begin{ex} %[5D4V22]
Một thửa ruộng hình chữ nhật có diện tích $161\text{ m}^2$ và chiều dài $17{,}5\text{ m}$. Chu vi của thửa ruộng đó là
\choice
{$26{,}7\text{ m}$}
{\True $53{,}4\text{ m}$}
{$48{,}6\text{ m}$}
{$52{,}4\text{ m}$}
\loigiai{
Chiều rộng thửa ruộng là:
$161 : 17{,}5 = 9{,}2\text{ (m)}$.\\
Chu vi thửa ruộng là:
$(17{,}5 + 9{,}2) \times 2 = 26{,}7 \times 2 = 53{,}4\text{ (m)}$.\\
\textbf{Đáp án của câu hỏi là B}
}
\end{ex}

\begin{ex} %[5D4V22]
Một thùng chứa $39\text{ l}$ mật ong. Người ta muốn rót đều toàn bộ số mật ong này vào các chai, mỗi chai chứa $0{,}75\text{ l}$. Số chai mật ong rót được là
\choice
{$48\text{ chai}$}
{\True $52\text{ chai}$}
{$56\text{ chai}$}
{$50\text{ chai}$}
\loigiai{
Số chai rót được là:
$39 : 0{,}75 = 52\text{ (chai)}$.\\
\textbf{Đáp án của câu hỏi là B}
}
\end{ex}

\begin{ex} %[5D4V22]
Giá trị của biểu thức $(15{,}6 + 24{,}9) : 2{,}7 - 3{,}5$ là
\choice
{$12$}
{\True $11{,}5$}
{$10{,}5$}
{$12{,}5$}
\loigiai{
Ta có:
$$(15{,}6 + 24{,}9) : 2{,}7 - 3{,}5 = 40{,}5 : 2{,}7 - 3{,}5 = 15 - 3{,}5 = 11{,}5.$$
\textbf{Đáp án của câu hỏi là B}
}
\end{ex}

\begin{ex} %[5D4V22]
Mua $3{,}5\text{ m}$ vải hoa hết $175\,000$ đồng. Mua $5\text{ m}$ vải trắng hết $200\,000$ đồng. Hỏi mỗi mét vải hoa đắt hơn mỗi mét vải trắng bao nhiêu tiền?
\choice
{$8\,000\text{ đồng}$}
{\True $10\,000\text{ đồng}$}
{$12\,000\text{ đồng}$}
{$15\,000\text{ đồng}$}
\loigiai{
Giá $1\text{ m}$ vải hoa là:
$175\,000 : 3{,}5 = 50\,000\text{ (đồng)}$.\\
Giá $1\text{ m}$ vải trắng là:
$200\,000 : 5 = 40\,000\text{ (đồng)}$.\\
Mỗi mét vải hoa đắt hơn vải trắng là:
$50\,000 - 40\,000 = 10\,000\text{ (đồng)}$.\\
\textbf{Đáp án của câu hỏi là B}
}
\end{ex}

\begin{ex} %[5D4V22]
Tìm số dư của phép chia $218 : 3{,}7$ nếu phần thập phân của thương lấy đến hai chữ số.
\choice
{$0{,}33$}
{$3{,}3$}
{\True $0{,}033$}
{$0{,}0033$}
\loigiai{
Thực hiện phép chia: $218 : 3{,}7 = 58{,}91$ (lấy đến hai chữ số ở phần thập phân).\\
Ta có: $58{,}91 \times 3{,}7 = 217{,}967$.\\
Số dư là: $218 - 217{,}967 = 0{,}033$.\\
\textbf{Đáp án của câu hỏi là C}
}
\end{ex}

\begin{ex} %[5D4V22]
Bác An rót $25{,}5\text{ l}$ nước mắm vào đầy các can loại $4{,}5\text{ l}$. Hỏi bác An cần ít nhất bao nhiêu chiếc can loại đó để chứa hết số nước mắm trên?
\choice
{$5\text{ can}$}
{\True $6\text{ can}$}
{$7\text{ can}$}
{$8\text{ can}$}
\loigiai{
Ta có: $25{,}5 : 4{,}5 = 5$ (dư $3\text{ l}$).\\
Như vậy rót được 5 can đầy và còn thừa $3\text{ l}$. Để chứa hết số nước mắm cần thêm 1 can nữa.\\
Số can ít nhất cần dùng là: $5 + 1 = 6\text{ (can)}$.\\
\textbf{Đáp án của câu hỏi là B}
}
\end{ex}

\begin{ex} %[5D4V22]
Một hình tam giác có diện tích $28{,}8\text{ cm}^2$ và độ dài đáy là $7{,}2\text{ cm}$. Chiều cao tương ứng với cạnh đáy đó là
\choice
{$4\text{ cm}$}
{\True $8\text{ cm}$}
{$6\text{ cm}$}
{$12\text{ cm}$}
\loigiai{
Chiều cao của hình tam giác là:
$28{,}8 \times 2 : 7{,}2 = 57{,}6 : 7{,}2 = 8\text{ (cm)}$.\\
\textbf{Đáp án của câu hỏi là B}
}
\end{ex}

\textbf{Phần 2. Trắc nghiệm đúng sai}
\setcounterpageref{ex}{0}

\begin{ex} %[5D4V22]
Cho hai phép chia: $A = 75 : 1{,}25$ và $B = 33{,}6 : 0{,}56$.
\choiceTF
{\True Thương của phép chia $A$ bằng $60$}
{Thương của phép chia $B$ bằng $6$}
{\True Thương của hai phép chia $A$ và $B$ bằng nhau}
{\True Khi cùng nhân số bị chia của cả hai phép chia trên với $10$ và giữ nguyên số chia thì thương mới của chúng đều bằng $600$}
\loigiai{
\begin{itemchoice}
\item Đúng. $A = 75 : 1{,}25 = 60$.
\item Sai. $B = 33{,}6 : 0{,}56 = 60 
eq 6$.
\item Đúng. Cả hai phép tính đều có thương là $60$.
\item Đúng. Khi nhân số bị chia với 10 và giữ nguyên số chia thì thương tăng 10 lần: $60 \times 10 = 600$.
\end{itemchoice}
}
\end{ex}

\begin{ex} %[5D4V22]
Một cửa hàng có $145\text{ kg}$ gạo nếp đóng vào các túi, mỗi túi chứa $3{,}5\text{ kg}$.
\choiceTF
{\True Phép chia $145 : 3{,}5$ có thương là $41$ và số dư là $1{,}5\text{ kg}$}
{Cửa hàng đóng được đúng $42$ túi gạo đầy}
{\True Lượng gạo còn dư chưa đủ đóng một túi là $1{,}5\text{ kg}$}
{\True Để đóng thêm vừa đủ 1 túi gạo nữa thì cửa hàng cần thêm ít nhất $2\text{ kg}$ gạo nếp}
\loigiai{
\begin{itemchoice}
\item Đúng. Ta có $145 : 3{,}5 = 41$ (dư $1{,}5$).
\item Sai. Chỉ đóng được nhiều nhất $41$ túi gạo đầy.
\item Đúng. Số gạo dư ra là $1{,}5\text{ kg}$.
\item Đúng. Số gạo cần thêm là: $3{,}5 - 1{,}5 = 2\text{ (kg)}$.
\end{itemchoice}
}
\end{ex}

\textbf{Phần 3. Trắc nghiệm trả lời ngắn}
\setcounterpageref{ex}{0}

\begin{ex} %[5D4V22]
Tính giá trị của biểu thức: $(72 : 0{,}5) : (18 \times 2)$.
\shortans{4}
\loigiai{
Ta có: $72 : 0{,}5 = 144$ và $18 \times 2 = 36$.\\
Giá trị biểu thức là: $144 : 36 = 4$.
}
\end{ex}

\begin{ex} %[5D4V22]
Một mảnh đất hình chữ nhật có diện tích $252\text{ m}^2$, chiều rộng đo được $10{,}5\text{ m}$. Chiều dài của mảnh đất đó bằng bao nhiêu mét?
\shortans{24}
\loigiai{
Chiều dài của mảnh đất là:
$252 : 10{,}5 = 24\text{ (m)}$.
}
\end{ex}

\begin{ex} %[5D4V22]
Tính tổng các số tự nhiên $x$ thỏa mãn: $12{,}5 < 2{,}5 \times x < 18$.
\shortans{13}
\loigiai{
Ta có: $12{,}5 : 2{,}5 < x < 18 : 2{,}5 \implies 5 < x < 7{,}2$.\\
Các số tự nhiên $x$ thỏa mãn là: $6$ và $7$.\\
Tổng của chúng là: $6 + 7 = 13$.
}
\end{ex}

\begin{ex} %[5D4V22]
Một thùng đựng $54\text{ l}$ dầu ăn. Người ta chia đều vào các can, mỗi can chứa $4{,}5\text{ l}$. Hỏi chia được tất cả bao nhiêu can như thế?
\shortans{12}
\loigiai{
Số can chia được là:
$54 : 4{,}5 = 12\text{ (can)}$.
}
\end{ex}

\begin{ex} %[5D4V22]
Người ta dùng một sợi dây thép dài $36\text{ m}$ để uốn thành các khung hình vuông, mỗi cạnh dài $0{,}75\text{ m}$. Hỏi uốn được nhiều nhất bao nhiêu khung hình vuông như thế?
\shortans{12}
\loigiai{
Chu vi một khung hình vuông là:
$0{,}75 \times 4 = 3\text{ (m)}$.\\
Số khung hình vuông uốn được là:
$36 : 3 = 12\text{ (khung)}$.
}
\end{ex}

\begin{ex} %[5D4V22]
Một xe máy chạy quãng đường dài $126\text{ km}$ hết $3{,}5$ giờ. Hỏi với vận tốc đó, xe máy chạy quãng đường $180\text{ km}$ thì hết thời gian bao nhiêu giờ?
\shortans{5}
\loigiai{
Vận tốc của xe máy là:
$126 : 3{,}5 = 36\text{ (km/h)}$.\\
Thời gian để đi quãng đường $180\text{ km}$ là:
$180 : 36 = 5\text{ (giờ)}$.
}
\end{ex}
